%=========== ΛΥΣΗ - 16ο ΘΕΜΑ Α ===========
\providecommand{\theoriaref}[3]{%
\begin{center}
\fcolorbox{maincolor!70!black}{maincolor!8}{%
\begin{minipage}{.88\linewidth}
\textcolor{maincolor!80!black}{\kerkissans{\textbf{#1~\ref{#2}}}}\\[1mm]
#3
\end{minipage}}
\end{center}}
\providecommand{\orismosref}[1]{%
\theoriaref{Ορισμός}{#1}{Η απάντηση δίνεται από τον αντίστοιχο ορισμό.}}
\providecommand{\apodeixiref}[1]{%
\theoriaref{Θεώρημα}{#1}{Η ζητούμενη απόδειξη δίνεται στην αντίστοιχη απόδειξη.}}

\begin{thema}{Α}
\begin{erwthma}
\item \theoriaref{Θεώρημα}{thewrhma:thmt}{Η διατύπωση και η γεωμετρική ερμηνεία δίνονται στην αντίστοιχη ενότητα.}

\item \orismosref{orismos:plagia_asymptwth}

\item \orismosref{orismos:krisima_shmeia}

\item Οι χαρακτηρισμοί των προτάσεων είναι:
\begin{alist}
\item \textbf{Λάθος}. Μια συνάρτηση μπορεί να παίρνει μέγιστη και ελάχιστη τιμή χωρίς να είναι συνεχής. Για παράδειγμα, η συνάρτηση
\[ f(x)=\begin{cases}0, & x<0\\ 1, & x\ge0\end{cases} \]
στο $[-1,1]$ παίρνει ελάχιστη και μέγιστη τιμή, αλλά δεν είναι συνεχής στο $0$.
\item \textbf{Σωστό}. Το πεδίο ορισμού της $f^{-1}$ είναι το σύνολο τιμών της $f$, δηλαδή το $f(A)$.
\item \textbf{Σωστό}. Αφού $f(0)<0$ και $\lim\limits_{x\to+\infty}f(x)>0$, υπάρχει $a>0$ με $f(a)>0$, άρα από το θεώρημα \eng{Bolzano} υπάρχει ρίζα στο $(0,a)$. Επειδή η $f$ είναι άρτια, υπάρχει και η συμμετρική της ρίζα στο $(-a,0)$.
\item \textbf{Σωστό}. Για κάθε σταθερά $c\in\mathbb{R}$ ισχύει
\[ \dint_a^\beta c\d x=c(\beta-a). \]
\item \textbf{Σωστό}. Από $0\le f(x)\le M$ κοντά στο $x_0$ έχουμε
\[ |(x-x_0)f(x)|\le M|x-x_0| \]
και επειδή $M|x-x_0|\to0$, από το κριτήριο παρεμβολής προκύπτει το ζητούμενο όριο.
\end{alist}

\item Σωστή απάντηση είναι η \textbf{β}. Αν $m\le f(x)\le M$ στο $[a,\beta]$, τότε
\[ m(\beta-a)\le\dint_a^\beta f(x)\d x\le M(\beta-a). \]
\end{erwthma}
\end{thema}
%=========== ΤΕΛΟΣ ΛΥΣΗΣ - 16ο ΘΕΜΑ Α ===========
